\subsection{Numerical Spectral Methods}

\begin{definition} \label{def:upper-hessenberg-matrix} \hlt{Upper Hessenberg Matrix}\\
Let $A=[a_{ij}]\in\operatorname{Mat}(n,n,\mathbb{F})$. Then $A$ is \hlt{upper Hessenberg} if $a_{ij}=0$ whenever $i>j+1$.\\
Equivalently, all entries below the first subdiagonal are zero.
\end{definition}

Numerical linear algebra also replaces a matrix by a unitarily similar one with many forced zero entries. Upper Hessenberg form is the standard preliminary reduction before QR-type eigenvalue algorithms.

\begin{figure}[H]
  \centering
  \begin{tikzpicture}[font=\small]
  \node at (0,0) {$
  \begin{bmatrix}
  * & * & * & * & *\\
  * & * & * & * & *\\
  0 & * & * & * & *\\
  0 & 0 & * & * & *\\
  0 & 0 & 0 & * & *
  \end{bmatrix}
  $};
  \draw[red!70, thick] (-1.35,-0.45) -- (0.65,-1.15);
  \node[red!70!black, align=center] at (2.45,-1.0) {first\\subdiagonal};
  \draw[->, red!70, thick] (1.65,-0.95) -- (0.72,-0.92);

  \node[blue!70!black, align=center] at (0,1.75) {entries below this subdiagonal are zero};
  \draw[->, blue!70, thick] (-0.2,1.45) -- (-1.0,-0.85);
\end{tikzpicture}

  \caption{Upper Hessenberg pattern: red is the first subdiagonal; blue entries below it are forced zero.}
  \label{fig:upper-hessenberg-pattern}
\end{figure}

\begin{theorem} \label{thm:unitary-upper-hessenberg-reduction} \hlt{Unitary Upper Hessenberg Reduction}\\
Every $A\in\operatorname{Mat}(n,n,\mathbb{F})$ is unitarily similar to upper Hessenberg matrix with real nonnegative subdiagonal entries.
\begin{proof}
Set $A_0=A$. For $k=1,\ldots,n-2$, suppose columns $1,\ldots,k-1$ of $A_{k-1}$ already have zeros below the first subdiagonal. Let $\mathbf{w}$ be the tail of column $k$ in rows $k+1,\ldots,n$.\\
By Corollary \ref{cor:householder-unitary-annihilator}, choose a unitary $V_k$ with $V_k^*\mathbf{w}=\alpha\mathbf{e}_1$, $\alpha\geq0$.\\
Put $U_k=I_k\oplus V_k$ and $A_k=U_k^*A_{k-1}U_k$. Right multiplication by $U_k$ changes only columns $k+1,\ldots,n$, so columns $1,\ldots,k-1$ keep their existing Hessenberg zeros. Left multiplication by $U_k^*$ acts only on rows $k+1,\ldots,n$; those rows are already zero in columns $1,\ldots,k-1$, so the zeros remain. In column $k$, the tail becomes $V_k^*\mathbf{w}=\alpha\mathbf{e}_1$, so entries below row $k+1$ are zero and the new subdiagonal entry is real nonnegative.\\
After $n-2$ steps, $A_{n-2}=U^*AU$ is upper Hessenberg, where $U=U_1\cdots U_{n-2}$ is unitary.\\
\hlt{Structural Commentary}: Each trailing unitary creates one Hessenberg column and preserves earlier ones.\\
\hlt{Key Insight}: Kill one column tail at a time.
\end{proof}
\end{theorem}

\begin{perspective} \label{per:hessenberg-before-qr} \hlt{Why Hessenberg Form Matters}\\
Hessenberg form preserves spectral data by unitary similarity, making QR-type eigenvalue algorithms cheaper.
\end{perspective}

\begin{method} \label{meth:qr-eigenvalue-algorithm-preview} \hlt{QR Eigenvalue Algorithm Preview}\\
For a square matrix $A_0$, the basic QR eigenvalue iteration is:
\begin{enumerate}[label=(\roman*)]
\setlength{\itemsep}{0pt}
\item In practice, first reduce $A_0$ to Hessenberg form by Theorem \ref{thm:unitary-upper-hessenberg-reduction}.
\item Factor $A_k=Q_kR_k$, with $Q_k$ unitary and $R_k$ upper triangular.
\item Set $A_{k+1}=R_kQ_k=Q_k^*A_kQ_k$, using shifts to speed convergence.
\end{enumerate}
Each step is a unitary similarity, so the spectrum is preserved.
\end{method}

\begin{method} \label{meth:hessenberg-qr-step-plane-rotations} \hlt{Hessenberg QR Step by Plane Rotations}\\
Let $H\in\operatorname{Mat}(n,n,\mathbb{C})$ be upper Hessenberg and let $\mu\in\mathbb{C}$.
\begin{enumerate}[label=(\roman*)]
\setlength{\itemsep}{0pt}
\item Set $B_0=H-\mu I_n$.
\item For $j=1,\ldots,n-1$, choose a unitary plane rotation $G_j$ acting on coordinates $j,j+1$,\\
so that $(G_j^*B_{j-1})_{j+1,j}=0$ and $B_j=G_j^*B_{j-1}$.
\item Put $Q=G_1\cdots G_{n-1}$ and $R=B_{n-1}$; then $H-\mu I_n=QR$.
\item Define the shifted QR update $H_+=RQ+\mu I_n=Q^*HQ$
\end{enumerate}
Adjacent rotations implement the Hessenberg step by local bulge chasing.
\end{method}

\begin{theorem} \label{thm:power-method-dominant-eigenvector-preview} \hlt{Power Method Dominant-Eigenvector Preview}\\
Let $A\mathbf{v}_j=\lambda_j\mathbf{v}_j$, where $\mathbf{v}_1,\ldots,\mathbf{v}_n$ is a basis of $\mathbb{C}^n$, and assume $\abs{\lambda_1}>\max_{j\geq2}\abs{\lambda_j}$.\\
If $\mathbf{u}_0=\sum\limits_{j=1}^{n}c_j\mathbf{v}_j$ with $c_1\neq0$, then $\lambda_1^{-k}A^k\mathbf{u}_0\longrightarrow c_1\mathbf{v}_1$. Thus the directions of $A^k\mathbf{u}_0$ converge to $\operatorname{span}\{\mathbf{v}_1\}$.
\begin{proof}
Set $\rho=\max\limits_{j\geq2}\abs{\lambda_j/\lambda_1}<1$. Since $A^k\mathbf{v}_j=\lambda_j^k\mathbf{v}_j$, $\lambda_1^{-k}A^k\mathbf{u}_0 =c_1\mathbf{v}_1+\sum\limits_{j=2}^{n}c_j\left(\frac{\lambda_j}{\lambda_1}\right)^k\mathbf{v}_j$.\\
The remainder has norm at most $\rho^k\sum\limits_{j=2}^{n}\abs{c_j}\norm{\mathbf{v}_j}$, hence tends to $0$.\\
\hlt{Structural Commentary}: The power method is spectral filtering by repeated multiplication.\\
\hlt{Key Insight}: A spectral gap makes every non-dominant component decay relative to the dominant one.
\end{proof}
\end{theorem}

\begin{method} \label{meth:inverse-and-shifted-power-iteration-preview} \hlt{Inverse and Shifted Power Iteration Preview}\\
Let $\mathcal{N}(\mathbf{z})\defeq\mathbf{z}/\norm{\mathbf{z}}_2$ for $\mathbf{z}\neq\mathbf{0}$.
\begin{enumerate}[label=(\roman*)]
\setlength{\itemsep}{0pt}
\item \hlt{Inverse}: if $0\notin\operatorname{spec}(A)$, solve $A\mathbf{y}_{k+1}=\mathbf{x}_k$ and set $\mathbf{x}_{k+1}=\mathcal{N}(\mathbf{y}_{k+1})$.\\
This is power iteration on $A^{-1}$ and targets minimal $\abs{\lambda}$.
\item \hlt{Shift-invert}: for $\sigma\notin\operatorname{spec}(A)$, $(A-\sigma I)\mathbf{y}_{k+1}=\mathbf{x}_k$, $\mathbf{x}_{k+1}=\mathcal{N}(\mathbf{y}_{k+1})$.\\
Eigenvalues become $(\lambda-\sigma)^{-1}$, targeting $\lambda$ nearest $\sigma$.
\item \hlt{Shifted power}: $\mathbf{x}_{k+1}=\mathcal{N}\bigl((A-\sigma I)\mathbf{x}_k\bigr)$; eigenvalues become $\lambda-\sigma$, emphasizing large $\abs{\lambda-\sigma}$.
\end{enumerate}
Use factors of $A$ or $A-\sigma I$ for solves; do not form explicit inverses.
\end{method}

\begin{method} \label{meth:svd-numerics-preview} \hlt{SVD Numerics Preview}\\
Stable numerical SVD computation proceeds as:
\begin{enumerate}[label=(\roman*)]
\setlength{\itemsep}{0pt}
\item Avoid forming $A^*A$ when possible, since it can square the condition number.
\item Use orthogonal or unitary transformations to reduce $A$ to bidiagonal form.
\item Compute singular values from the bidiagonal problem, then use them for low-rank approximation, pseudoinverse truncation, and least-squares diagnostics through Theorem \ref{thm:eckart-young-frobenius-spectral} and Definition \ref{def:spectral-condition-number}.
\end{enumerate}
\end{method}

\begin{definition} \label{def:discrete-fourier-matrix} \hlt{Discrete Fourier Matrix}\\
For $n\geq1$, set $\omega_n=e^{2\pi i/n}$ and $F_n\defeq[\omega_n^{jk}]_{j,k=0}^{n-1}$.\\
Then $F_n\mathbf{c}$ is the unnormalised DFT of $\mathbf{c}\in\C^n$; $Q_n\defeq n^{-1/2}F_n$ is the normalized Fourier matrix.
\end{definition}

\begin{definition} \label{def:circulant-matrix-cyclic-shift} \hlt{Circulant Matrix}\\
On the $0$-indexed basis of $\C^n$, define the \hlt{cyclic shift}
$P_n\mathbf{e}_j=\mathbf{e}_{(j-1)\operatorname{mod} n}$.\\
A matrix $C\in\operatorname{Mat}(n,n,\C)$ is \hlt{circulant} if, for some $a_0,\ldots,a_{n-1}\in\C$,
\begin{align}
C&=a_0I_n+a_1P_n+\cdots+a_{n-1}P_n^{n-1}=p(P_n), \nonumber\\
p(z)&=a_0+a_1z+\cdots+a_{n-1}z^{n-1} \nonumber
\end{align}
\end{definition}

\begin{theorem} \label{thm:fourier-basis-diagonalizes-cyclic-circulants} \hlt{Fourier Basis for Cyclic Shifts and Circulants}\\
Let $n\geq1$, let $\omega_n\defeq e^{2\pi i/n}$, and let $P_n$ be the cyclic shift.\\
For $0\leq k<n$, set $\mathbf{f}_k\defeq(\omega_n^{jk})_{j=0}^{n-1}$; $F_n\defeq[\mathbf{f}_0\ \cdots\ \mathbf{f}_{n-1}]$; and $\Omega_n\defeq\operatorname{diag}(1,\omega_n,\ldots,\omega_n^{n-1})$.\\
Then $F_n^*F_n=nI_n$, so $Q_n\defeq n^{-1/2}F_n$ is unitary, and $P_nF_n=F_n\Omega_n$.\\
If $C=p(P_n)$ is circulant, then $Q_n^*CQ_n=p(\Omega_n) =\operatorname{diag}\bigl(p(1),p(\omega_n),\ldots,p(\omega_n^{n-1})\bigr)$.
\begin{proof}
For $0\leq k,\ell\leq n-1$, $\innerprod{\mathbf{f}_k,\mathbf{f}_\ell} =\sum\limits_{r=0}^{n-1}(\omega_n^{k-\ell})^r =n\delta_{k\ell}$.
For $k=\ell$ the sum is $n$; for $k\neq\ell$, $W=\omega_n^{k-\ell}$ has $W\neq1$ and $W^n=1$, so $(1-W)\sum\limits_{r=0}^{n-1}W^r=0$. Hence $F_n^*F_n=nI_n$, and Theorem \ref{thm:unitary-matrix-equivalent-conditions} makes $Q_n$ unitary.\\
Also $P_n\mathbf{f}_k=\omega_n^k\mathbf{f}_k$, so $Q_n^*P_nQ_n=\Omega_n$; applying $p$ gives the formula.\\
\hlt{Structural Commentary}: Fourier coordinates diagonalize every circulant.\\
\hlt{Key Insight}: Orthogonality of roots of unity turns a cyclic matrix problem into scalar evaluation.
\end{proof}
\end{theorem}

\begin{figure}[H]
  \centering
  \begin{tikzpicture}[
  >=Stealth,
  x=0.92cm,
  y=0.9cm,
  every node/.style={font=\small},
  box/.style={draw=gray!55, rounded corners=3pt, fill=gray!6, align=center, minimum width=2.15cm, minimum height=0.8cm},
  result/.style={box, fill=green!7},
  arrow/.style={->, line width=0.85pt, blue!70!black},
  lbl/.style={fill=white, text=gray!70!black, inner xsep=2pt, inner ysep=1pt}
]
  \node[box] (shift) at (0,1.1) {$P_n$};
  \node[box] (diagshift) at (4.15,1.1) {$\Omega_n$};
  \node[box] (circ) at (0,-1.1) {$C=p(P_n)$};
  \node[result] (diagcirc) at (4.15,-1.1) {$p(\Omega_n)$};

  \draw[arrow] (shift) -- node[lbl, above] {$Q_n^*(\,\cdot\,)Q_n$} (diagshift);
  \draw[arrow] (circ) -- node[lbl, below] {$Q_n^*(\,\cdot\,)Q_n$} (diagcirc);
  \draw[arrow] (shift) -- node[lbl, left] {$p(\cdot)$} (circ);
  \draw[arrow] (diagshift) -- node[lbl, right] {$p(\cdot)$} (diagcirc);
\end{tikzpicture}

  \caption{Fourier diagonalizes $P_n$; polynomial evaluation diagonalizes $C=p(P_n)$.}
  \label{fig:circulant-fourier-diagonalization}
\end{figure}

\begin{proposition} \label{prop:fourier-matrix-inverse} \hlt{Inverse of the Fourier Matrix}\\
The Fourier matrix satisfies $F_n^*F_n=nI_n$, and therefore $F_n^{-1}=\frac1nF_n^*$.
\begin{proof}
Theorem \ref{thm:fourier-basis-diagonalizes-cyclic-circulants} gives $Q_n^*Q_n=I_n$ for $Q_n=n^{-1/2}F_n$. Multiplying by $n$ gives $F_n^*F_n=nI_n$, so $n^{-1}F_n^*$ is a left inverse of $F_n$. Since $F_n$ is square, Theorem \ref{thm:square-one-sided-inverse-is-two-sided} makes it the inverse.\\
\hlt{Structural Commentary}: Fourier inversion is conjugate transpose plus one scalar normalization.\\
\hlt{Key Insight}: The inverse is simple because the Fourier columns are orthogonal.
\end{proof}
\end{proposition}

\begin{method} \label{meth:radix-two-fft-split} \hlt{Radix-Two FFT Split}\\
Let $n=2m$ and $\mathbf{y}=F_n\mathbf{c}$.
\begin{enumerate}[label=(\roman*)]
\setlength{\itemsep}{0pt}
\item Split $\mathbf{c}'\defeq(c_{2r})_{r=0}^{m-1}$ and $\mathbf{c}''\defeq(c_{2r+1})_{r=0}^{m-1}$.
\item Compute $\mathbf{y}'\defeq F_m\mathbf{c}'$ and $\mathbf{y}''\defeq F_m\mathbf{c}''$.
\item With $d_j\defeq\omega_n^j$, and using $\omega_n^2=\omega_m$ and $\omega_n^m=-1$, set $y_j=y'_j+d_jy''_j$; and $y_{j+m}=y'_j-d_jy''_j$ for $0\leq j<m$.
\item For $n=2^\ell$, repeat (i)--(iii) for $\ell$ stages.\\
$n^2$ dense products become about $\frac12 n\log_2 n$ twiddle multiplications plus additions.
\end{enumerate}
\end{method}

\begin{figure}[H]
  \centering
  \begin{tikzpicture}[
  >=Stealth,
  every node/.style={font=\small},
  vector/.style={draw=gray!60, fill=white, rounded corners=2pt, minimum width=1.25cm, minimum height=0.58cm, align=center},
  block/.style={draw=gray!65, fill=gray!8, rounded corners=3pt, minimum width=1.25cm, minimum height=0.62cm, align=center},
  map/.style={->, line width=0.9pt, blue!70!black}
]
  \node[vector] (c) at (0,0) {$\mathbf{c}$};
  \node[block] (split) at (1.85,0) {split};

  \node[vector] (ce) at (3.75,0.8) {$\mathbf{c}'=(c_0,c_2)^t$};
  \node[vector] (co) at (3.75,-0.8) {$\mathbf{c}''=(c_1,c_3)^t$};

  \node[block] (f2e) at (5.8,0.8) {$F_2$};
  \node[block] (f2o) at (5.8,-0.8) {$F_2$};

  \node[vector] (ye) at (7.55,0.8) {$\mathbf{y}'$};
  \node[vector] (yo) at (7.55,-0.8) {$\mathbf{y}''$};

  \node[block, minimum width=1.45cm, minimum height=0.82cm] (butterfly) at (9.75,0) {$B_4$};
  \node[vector] (y) at (11.85,0) {$\mathbf{y}=F_4\mathbf{c}$};

  \draw[map] (c) -- (split);
  \draw[map] (split.east) -- (ce.west);
  \draw[map] (split.east) -- (co.west);
  \draw[map] (ce) -- (f2e);
  \draw[map] (co) -- (f2o);
  \draw[map] (f2e) -- (ye);
  \draw[map] (f2o) -- (yo);
  \draw[map] (ye.east) -- (butterfly.west);
  \draw[map] (yo.east) -- (butterfly.west);
  \draw[map] (butterfly) -- (y);

\end{tikzpicture}

  \caption{For $n=4$: split, apply two $F_2$ transforms, then $B_4\binom{\mathbf{y}'}{\mathbf{y}''}=\binom{\mathbf{y}'+D\mathbf{y}''}{\mathbf{y}'-D\mathbf{y}''}$ with $D=\operatorname{diag}(1,i)$.}
  \label{fig:fft-butterfly-n4}
\end{figure}
